% Options for packages loaded elsewhere
\PassOptionsToPackage{unicode}{hyperref}
\PassOptionsToPackage{hyphens}{url}
\documentclass[
  ignorenonframetext,
]{beamer}
\newif\ifbibliography
\usepackage{pgfpages}
\setbeamertemplate{caption}[numbered]
\setbeamertemplate{caption label separator}{: }
\setbeamercolor{caption name}{fg=normal text.fg}
\beamertemplatenavigationsymbolsempty
% remove section numbering
\setbeamertemplate{part page}{
  \centering
  \begin{beamercolorbox}[sep=16pt,center]{part title}
    \usebeamerfont{part title}\insertpart\par
  \end{beamercolorbox}
}
\setbeamertemplate{section page}{
  \centering
  \begin{beamercolorbox}[sep=12pt,center]{section title}
    \usebeamerfont{section title}\insertsection\par
  \end{beamercolorbox}
}
\setbeamertemplate{subsection page}{
  \centering
  \begin{beamercolorbox}[sep=8pt,center]{subsection title}
    \usebeamerfont{subsection title}\insertsubsection\par
  \end{beamercolorbox}
}
% Prevent slide breaks in the middle of a paragraph
\widowpenalties 1 10000
\raggedbottom
\AtBeginPart{
  \frame{\partpage}
}
\AtBeginSection{
  \ifbibliography
  \else
    \frame{\sectionpage}
  \fi
}
\AtBeginSubsection{
  \frame{\subsectionpage}
}
\usepackage{iftex}
\ifPDFTeX
  \usepackage[T1]{fontenc}
  \usepackage[utf8]{inputenc}
  \usepackage{textcomp} % provide euro and other symbols
\else % if luatex or xetex
  \usepackage{unicode-math} % this also loads fontspec
  \defaultfontfeatures{Scale=MatchLowercase}
  \defaultfontfeatures[\rmfamily]{Ligatures=TeX,Scale=1}
\fi
\usepackage{lmodern}
\ifPDFTeX\else
  % xetex/luatex font selection
\fi
% Use upquote if available, for straight quotes in verbatim environments
\IfFileExists{upquote.sty}{\usepackage{upquote}}{}
\IfFileExists{microtype.sty}{% use microtype if available
  \usepackage[]{microtype}
  \UseMicrotypeSet[protrusion]{basicmath} % disable protrusion for tt fonts
}{}
\makeatletter
\@ifundefined{KOMAClassName}{% if non-KOMA class
  \IfFileExists{parskip.sty}{%
    \usepackage{parskip}
  }{% else
    \setlength{\parindent}{0pt}
    \setlength{\parskip}{6pt plus 2pt minus 1pt}}
}{% if KOMA class
  \KOMAoptions{parskip=half}}
\makeatother
\setlength{\emergencystretch}{3em} % prevent overfull lines
\providecommand{\tightlist}{%
  \setlength{\itemsep}{0pt}\setlength{\parskip}{0pt}}
% Slide layout defaults for warning-clean scientific decks.
\usepackage{etoolbox}
\IfFileExists{xurl.sty}{\usepackage{xurl}}{}
\IfFileExists{seqsplit.sty}{\usepackage{seqsplit}}{\newcommand{\seqsplit}[1]{#1}}
\protected\def\breakseq#1{\seqsplit{#1}}
\protected\def\breaktt#1{\begingroup\ttfamily\seqsplit{#1}\endgroup}
\setlength{\emergencystretch}{6em}
\tolerance=5000
\hbadness=10000
\hfuzz=1pt
\setlength{\tabcolsep}{2pt}
\AtBeginEnvironment{longtable}{\tiny\renewcommand{\arraystretch}{0.86}\setlength{\tabcolsep}{1pt}}
\AtBeginEnvironment{tabular}{\tiny\renewcommand{\arraystretch}{0.86}\setlength{\tabcolsep}{1pt}}

% Natbib and cross-reference fallbacks — slides don't load natbib
% or cleveref, but combined-PDF manuscript prose may emit these
% commands. The fallback renders citations as a bracketed key list
% and cross-references as detokenized labels so slides don't fail on
% undefined control sequences or raw underscores. \providecommand is
% a no-op if packages load later.
\providecommand{\citep}[1]{[#1]}
\providecommand{\citet}[1]{#1}
\providecommand{\citealp}[1]{#1}
\providecommand{\citeauthor}[1]{#1}
\providecommand{\citeyear}[1]{#1}
\providecommand{\cref}[1]{\texttt{\detokenize{#1}}}
\providecommand{\Cref}[1]{\texttt{\detokenize{#1}}}
\usepackage{bookmark}
\IfFileExists{xurl.sty}{\usepackage{xurl}}{} % add URL line breaks if available
\urlstyle{same}
% fallback for those not using the hyperref driver hyperxmp:
\makeatletter
\@ifundefined{xmpquote}{\newcommand{\xmpquote}[1]{#1}}{}
\makeatother
\hypersetup{
  hidelinks,
  pdfcreator={LaTeX via pandoc}}

\author{\texorpdfstring{}{}}
\date{}

\begin{document}

\section{Abstract}\label{abstract}

\begin{frame}[fragile,allowframebreaks]{Abstract}
The reproducibility crisis in computational research is fundamentally
structural: research artifacts are scattered across disconnected
tools---LaTeX editors, Jupyter notebooks, ad-hoc shell scripts---with no
enforced mechanism to keep code, data, and manuscript synchronized.
Studies have shown that most published findings are false positives,
replication rates in psychology hover around 36\%, and only 24\% of 1.4
million Jupyter notebooks can be successfully re-executed. Existing
tools address fragments of this problem: workflow managers (Snakemake,
Nextflow, CWL) orchestrate computation; literate programming systems
(Quarto, Jupyter Book, R Markdown, Overleaf, OpenAI Prism) render
documents; data versioning tools (DVC) track artifacts---but none
enforces cross-cutting quality standards as architectural invariants.
\texttt{template/} applies the principle of Infrastructure as Code to
the research lifecycle, making the manuscript, test suite, and
provenance chain version-controlled, deterministically buildable, and
independently verifiable. It is built on a Two-Layer Architecture that
separates 23 infrastructure subdirectories (20 importable Python
packages, \textasciitilde567 modules, validated by \textasciitilde7,375
tests) from self-contained project workspaces, connected by a
YAML-declared pipeline (12 stages; default full 10)-based build pipeline
progressing from environment sanitization through test execution (with a
Zero-Mock testing policy enforcing 90\% project-level and 60\%
infrastructure-level coverage via real filesystem operations and
subprocess invocations), analysis script invocation, Pandoc/XeLaTeX
rendering, SHA-256 cryptographic hashing with steganographic
watermarking, structural PDF validation, and LLM-assisted review. A
Documentation Duality standard equips every directory with both
human-readable \texttt{README.md} and machine-readable
\texttt{AGENTS.md} files, while each infrastructure module additionally
carries a \texttt{SKILL.md}---a structured skill descriptor aligned with
the Model Context Protocol---enabling AI agents to locate and invoke
module capabilities without hallucinating API signatures.

Scalability is demonstrated across the generated public exemplar roster
(\breaktt{templates/template\_active\_inference},
\breaktt{templates/template\_autoresearch\_project},
\breaktt{templates/template\_autoscientists},
\breaktt{templates/template\_code\_project},
\breaktt{templates/template\_madlib},
\breaktt{templates/template\_newspaper},
\breaktt{templates/template\_prose\_project},
\breaktt{templates/template\_sia}, \breaktt{templates/template\_template},
\breaktt{templates/template\_textbook}), with representative
heterogeneous cases under \breaktt{projects/templates/}: optimization
(\breaktt{template\_code\_project}, 197 tests), prose
(\breaktt{template\_prose\_project}, 78 tests), and AutoResearch
readiness (\breaktt{template\_autoresearch\_project}, 157 tests). These
guarantee control-positive layouts for code-centric, prose-centric, and
retrieval-centric workflows at 90\%+ project coverage alongside 60\%+
infrastructure gates. All three share identical pipeline stages without
cross-project coupling. This manuscript adds a complementary reflexive
artifact: authored from \breaktt{projects/templates/template\_template}
(89 tests) as a public exemplar in the same discovered/rendered tree,
using the same analysis and render path and injecting counters from
repository introspection. The fact that these words, metrics, and
figures were generated by the pipeline they describe demonstrates
self-documenting capacity: rendered through the DAG, validated without
mocks, optionally watermarked. A comparative analysis against nine peer
tools across fourteen dimensions positions \texttt{template/} as
integrating fourteen distinctive enforcement capabilities---testing
thresholds, cryptographic provenance, steganographic watermarking,
multi-project management, MCP-aligned skill descriptors, Zero-Mock
policy, orchestration through publication---in one repository. Code is
released under the Apache License 2.0 at
\breaktt{github.com/docxology/template}; the work remains open-ended.
\end{frame}

\end{document}
